
\documentclass{article}
\usepackage{colortbl}
\usepackage{makecell}
\usepackage{multirow}
\usepackage{supertabular}

\begin{document}

\newcounter{utterance}

\centering \large Interaction Transcript for game `matchit\_ascii', experiment `similar\_grid\_1', episode 7 with reka{-}flash{-}3.1{-}t0.0.
\vspace{24pt}

{ \footnotesize  \setcounter{utterance}{1}
\setlength{\tabcolsep}{0pt}
\begin{supertabular}{c@{$\;$}|p{.15\linewidth}@{}p{.15\linewidth}p{.15\linewidth}p{.15\linewidth}p{.15\linewidth}p{.15\linewidth}}
    \# & $\;$A & \multicolumn{4}{c}{Game Master} & $\;\:$B\\
    \hline

    \theutterance \stepcounter{utterance}  
    & & \multicolumn{4}{p{0.6\linewidth}}{
        \cellcolor[rgb]{0.9,0.9,0.9}{
            \makecell[{{p{\linewidth}}}]{
                \texttt{\tiny{[P1$\langle$GM]}}
                \texttt{You are participating in a collaborative guessing game.} \\
\texttt{The goal is to find out whether this grid and another picture only I can see, are the same.} \\
\texttt{Please describe your grid first. Then, I will provide my description and we can ask each other questions about the grids to figure out whether they are the same. Now start your short grid description with "DESCRIPTION:" followed by the description. Do not add anything else. This is your grid:} \\
\texttt{▢ ▢ ▢ ▢ X} \\
\texttt{▢ ▢ ▢ X ▢} \\
\texttt{▢ ▢ X ▢ ▢} \\
\texttt{▢ X ▢ ▢ ▢} \\
\texttt{X ▢ ▢ ▢ ▢} \\
            }
        }
    }
    & \\ \\

    \theutterance \stepcounter{utterance}  
    & \multicolumn{4}{p{0.6\linewidth}}{
        \cellcolor[rgb]{0.9,0.9,0.9}{
            \makecell[{{p{\linewidth}}}]{
                \texttt{\tiny{[P1$\rangle$GM]}}
                \texttt{<reasoning>} \\
\texttt{Let me try to figure out how to approach this. The user wants me to describe my grid so they can compare it with their own. The grid they provided has five rows, each with five characters. The characters are either ▢ (which I assume stands for empty) or X (which is probably the filled spot). Let me look at the structure again.} \\
\\ 
\texttt{First row: ▢ ▢ ▢ ▢ X. So the first four are empty, and the fifth is X. That would be the rightmost position in the first row.} \\
\\ 
\texttt{Second row: ▢ ▢ ▢ X ▢. Here, the fourth position is X, others empty. So that's one position to the left of the rightmost in the second row.} \\
\\ 
\texttt{Third row: ▢ ▢ X ▢ ▢. The third position is X. So centered in the third row.} \\
\\ 
\texttt{Fourth row: ▢ X ▢ ▢ ▢. The second position is X. So one to the right of the leftmost in the fourth row.} \\
\\ 
\texttt{Fifth row: X ▢ ▢ ▢ ▢. The first position is X. Leftmost in the fifth row.} \\
\\ 
\texttt{Wait, but the user's grid is written as:} \\
\\ 
\texttt{▢ ▢ ▢ ▢ X} \\
\\ 
\texttt{▢ ▢ ▢ X ▢} \\
\\ 
\texttt{▢ ▢ X ▢ ▢} \\
\\ 
\texttt{▢ X ▢ ▢ ▢} \\
\\ 
\texttt{X ▢ ▢ ▢ ▢} \\
\\ 
\texttt{So each row is five elements. Let me visualize each row's positions as columns 1 to 5.} \\
\\ 
\texttt{Row 1: Columns 1{-}4 are empty, column 5 is X.} \\
\\ 
\texttt{Row 2: Columns 1{-}3 empty, column 4 X, column 5 empty.} \\
\\ 
\texttt{Row 3: Columns 1{-}2 empty, column 3 X, columns 4{-}5 empty.} \\
\\ 
\texttt{Row 4: Columns 1 empty, column 2 X, columns 3{-}5 empty.} \\
\\ 
\texttt{Row 5: Column 1 X, columns 2{-}5 empty.} \\
\\ 
\texttt{So the X positions are:} \\
\\ 
\texttt{Row 1, Column 5} \\
\\ 
\texttt{Row 2, Column 4} \\
\\ 
\texttt{Row 3, Column 3} \\
\\ 
\texttt{Row 4, Column 2} \\
\\ 
\texttt{Row 5, Column 1} \\
\\ 
\texttt{Wait, but that seems like each X is in a different row and column. Let me check again.} \\
\\ 
\texttt{Row 1: X at position 5} \\
            }
        }
    }
    & & \\ \\

    \theutterance \stepcounter{utterance}  
    & & & \multicolumn{2}{p{0.3\linewidth}}{
        \cellcolor[rgb]{0.9,0.9,0.9}{
            \makecell[{{p{\linewidth}}}]{
                \texttt{\tiny{[GM$|$GM]}}
                \texttt{abort, first word: <reasoning>, but expected DESCRIPTION:} \\
            }
        }
    }
    & & \\ \\

    \theutterance \stepcounter{utterance}  
    & & & \multicolumn{2}{p{0.3\linewidth}}{
        \cellcolor[rgb]{0.9,0.9,0.9}{
            \makecell[{{p{\linewidth}}}]{
                \texttt{\tiny{[GM$|$GM]}}
                \texttt{Aborted} \\
            }
        }
    }
    & & \\ \\

\end{supertabular}
}

\end{document}
